\documentclass[11pt,a4paper]{article}

% ---------------------------------------------------------------------------
% Pacotes
% ---------------------------------------------------------------------------
\usepackage[utf8]{inputenc}
\usepackage[T1]{fontenc}
\usepackage[brazil]{babel}
\usepackage{lmodern}
\usepackage{microtype}
\usepackage[margin=2.5cm]{geometry}
\usepackage{amsmath,amssymb}
\usepackage{booktabs}
\usepackage{tabularx}
\usepackage{array}
\usepackage{enumitem}
\usepackage[table]{xcolor}
\usepackage[most]{tcolorbox}
\usepackage[hidelinks]{hyperref}

\setlist{itemsep=2pt,topsep=4pt}
\renewcommand{\arraystretch}{1.25}
\newcolumntype{L}{>{\raggedright\arraybackslash}X}

\definecolor{prodcolor}{RGB}{28,86,140}
\definecolor{notecolor}{RGB}{110,110,110}

% Caixa para a visão de produto
\newtcolorbox{produto}[1]{
  enhanced, breakable,
  colback=prodcolor!4, colframe=prodcolor,
  fonttitle=\bfseries, title={Visão de produto: #1},
  left=8pt, right=8pt, top=6pt, bottom=6pt
}

% Caixa para observações
\newtcolorbox{nota}{
  enhanced, breakable,
  colback=notecolor!6, colframe=notecolor!60,
  left=6pt, right=6pt, top=4pt, bottom=4pt
}

\newcommand{\prodsec}[1]{\par\medskip\noindent\textbf{\color{prodcolor}#1}\par\smallskip}

\title{Aprendizado de Máquina e Aprendizado Profundo\\no Setor Energético\\[4pt]
\large Aplicações, fundamentos técnicos e visão de produto}
\author{Mapeamento Energético}
\date{Setembro de 2026}

\begin{document}

\maketitle

\begin{abstract}
Este documento aprofunda quatro frentes de aplicação de aprendizado de máquina
(ML) e aprendizado profundo (DL) no setor elétrico: (1)~previsão de carga,
geração renovável e preço; (2)~operação e otimização da rede;
(3)~detecção de falhas, proteção e manutenção preditiva; e (4)~distribuição e
consumidor final. Para cada frente são apresentados o problema de negócio, a
formulação matemática, as famílias de modelos, os dados necessários, as
métricas de avaliação e os desafios práticos. Cada seção termina com
uma visão de produto: quem é o usuário, que problema se resolve, qual o
escopo mínimo viável, a arquitetura de referência, as métricas de sucesso, o
modelo de negócio e os riscos. O contexto regulatório e operacional brasileiro
(ONS, CCEE, ANEEL, distribuidoras) é usado como pano de fundo sempre que
relevante.
\end{abstract}

\tableofcontents
\newpage

% ===========================================================================
\section{Introdução}
% ===========================================================================

O setor elétrico passa por três transformações simultâneas que ampliam o
espaço para ML/DL:

\begin{enumerate}
  \item \textbf{Descarbonização}: a participação crescente de fontes eólica e
    solar, que são variáveis e não despacháveis, desloca o problema central
    da operação de ``atender a uma carga previsível com geração controlável''
    para ``equilibrar incertezas dos dois lados''.
  \item \textbf{Descentralização}: a micro e minigeração distribuída (MMGD),
    as baterias e os veículos elétricos transformam o consumidor em agente
    ativo e tornam a rede de distribuição bidirecional.
  \item \textbf{Digitalização}: medidores inteligentes, unidades de medição
    fasorial (PMUs), sensores em ativos, drones e imagens de satélite geram
    volumes de dados que métodos clássicos não exploram bem.
\end{enumerate}

Apesar do grande volume de pesquisa, os reviews recentes convergem em um
diagnóstico: há muitos pilotos e poucos sistemas em escala, e a barreira
principal não é a acurácia do modelo, mas a integração com sistemas legados, a
confiabilidade e a aceitação pelos operadores~\cite{energyinformatics2025,trustworthy2024}.
Por isso, cada seção deste documento trata o modelo e o produto
como problemas igualmente importantes.

\begin{nota}
\textbf{Como ler a visão de produto.} As visões propostas são hipóteses de
produto, não planos validados. Servem para orientar a descoberta: cada
hipótese de valor deve ser testada com usuários reais antes de virar escopo.
\end{nota}

% ===========================================================================
\section{Previsão: carga, geração renovável e preço}
% ===========================================================================

\subsection{Contexto e importância}

Previsão é a aplicação mais madura de ML no setor. Praticamente toda decisão
operacional e comercial depende de uma estimativa do futuro:

\begin{itemize}
  \item o \textbf{operador do sistema} (no Brasil, o ONS) precisa prever a carga
    e a geração renovável para programar o despacho, dimensionar reservas e
    avaliar a segurança;
  \item \textbf{geradores e comercializadores} precisam prever produção e preço
    para contratar, fazer \textit{hedge} e evitar exposição ao mercado de curto
    prazo;
  \item \textbf{distribuidoras} precisam prever a demanda para compra de energia
    em leilões, planejamento da expansão e gestão de carregamento de
    transformadores.
\end{itemize}

O caso mais citado de valor gerado é o da DeepMind com o Google: uma rede
neural treinada com previsões meteorológicas e histórico de turbinas passou
a prever a geração de 700\,MW de parques eólicos com 36 horas de
antecedência. Com isso, o Google passou a oferecer a energia no mercado
do dia seguinte, e a empresa reporta aumento de cerca de 20\% no valor da
energia eólica~\cite{deepmindwind}.

\subsection{Horizontes de previsão}

\begin{table}[h]
\centering
\begin{tabularx}{\textwidth}{l l L}
\toprule
\textbf{Horizonte} & \textbf{Faixa típica} & \textbf{Uso principal} \\
\midrule
Muito curto prazo & segundos a 1\,h & controle de frequência, rampas,
  despacho intra-horário, controle de inversores \\
Curto prazo & 1\,h a 7 dias & programação diária da operação, mercado do dia
  seguinte, unit commitment \\
Médio prazo & semanas a meses & planejamento da operação (PMO), manutenção,
  gestão de reservatórios, contratação \\
Longo prazo & anos & expansão da geração e da transmissão, leilões,
  planejamento tarifário \\
\bottomrule
\end{tabularx}
\caption{Horizontes de previsão e seus usos.}
\end{table}

Cada horizonte tem dados, modelos e métricas diferentes. Nos horizontes muito
curtos predominam dados locais de alta resolução (medições, câmeras de céu,
imagens de satélite). No curto prazo, as previsões numéricas de tempo (NWP)
são a principal variável explicativa. No médio e longo prazo, entram
cenários macroeconômicos e climáticos, e o problema se aproxima mais de
modelagem estrutural do que de ML puro.

\subsection{Formulação do problema}

Seja $y_t$ a variável-alvo (carga, potência gerada ou preço) no instante $t$,
$\mathbf{x}_t$ um vetor de covariáveis observadas (temperatura, irradiância
medida etc.) e $\mathbf{z}_{t+h}$ covariáveis conhecidas ou previstas para o
futuro (calendário, feriados, NWP). A previsão pontual com horizonte $h$ e
janela de contexto $L$ é
\begin{equation}
  \hat{y}_{t+h} = f_\theta\big(y_{t-L+1:t},\; \mathbf{x}_{t-L+1:t},\;
  \mathbf{z}_{t+1:t+h}\big).
\end{equation}

Na prática, a decisão raramente depende só do valor esperado. Dimensionar
reserva, fazer oferta em mercado ou decidir se uma térmica deve ser ligada
exige conhecer a \emph{distribuição} do futuro. Por isso a tendência é de
\textbf{previsão probabilística}, geralmente por quantis $q_\tau$ treinados
com a \textit{pinball loss}:
\begin{equation}
  \mathcal{L}_\tau(y, q_\tau) =
  \max\big(\tau\,(y - q_\tau),\; (\tau - 1)\,(y - q_\tau)\big),
  \qquad \tau \in (0,1).
\end{equation}
A qualidade da distribuição completa é medida pelo \textit{Continuous Ranked
Probability Score} (CRPS), em que $F$ é a distribuição prevista:
\begin{equation}
  \mathrm{CRPS}(F, y) = \int_{-\infty}^{\infty}
  \big(F(z) - \mathbb{1}\{z \ge y\}\big)^2 \, dz.
\end{equation}

\subsection{Previsão de carga}

\paragraph{Determinantes da carga.} A carga elétrica agregada tem padrões
fortes e relativamente estáveis:
\begin{itemize}
  \item \textbf{sazonalidade múltipla}: diária, semanal e anual;
  \item \textbf{efeitos de calendário}: feriados, pontes, horário de verão
    (quando existe), grandes eventos;
  \item \textbf{clima}: temperatura é a principal variável (relação não
    linear, em ``U'', por causa de ar-condicionado e aquecimento), seguida de
    umidade, sensação térmica e nebulosidade;
  \item \textbf{economia}: atividade industrial e comercial, choques como a
    pandemia de 2020.
\end{itemize}

\paragraph{Carga bruta versus carga líquida.} Com a expansão da MMGD solar, o
que o operador enxerga nos pontos de conexão é a \emph{carga líquida}: o
consumo menos a geração atrás do medidor. Em dias ensolarados, a carga
líquida afunda ao meio-dia e sobe de forma abrupta no fim da tarde, formando
a chamada ``curva do pato''. Prever a carga líquida diretamente mistura dois
fenômenos com dinâmicas diferentes. Por isso, a boa prática é decompor:
\begin{equation}
  \text{carga líquida}_t = \text{carga bruta}_t - \text{MMGD}_t,
\end{equation}
prevendo cada termo com suas próprias covariáveis. O ONS mantém um sistema
específico de previsão da geração da MMGD fotovoltaica, com horizonte de
72 horas, justamente por esse motivo~\cite{onssolar}.

\paragraph{Previsão hierárquica.} A carga existe em vários níveis
(subsistema, área, subestação, alimentador, transformador, consumidor).
Previsões independentes em cada nível não somam de forma coerente. Técnicas de
\textbf{reconciliação hierárquica} (como MinT) ajustam as previsões para que
os níveis sejam consistentes entre si e costumam melhorar a acurácia em todos
os níveis.

\subsection{Previsão de geração eólica}

A potência extraída do vento cresce com o cubo da velocidade:
\begin{equation}
  P = \tfrac{1}{2}\,\rho\, A\, C_p(\lambda, \beta)\, v^3,
\end{equation}
em que $\rho$ é a densidade do ar, $A$ a área varrida, $C_p$ o coeficiente de
potência e $v$ a velocidade do vento. Por causa do cubo, um erro pequeno na
previsão do vento vira um erro grande na previsão de potência, principalmente
na região íngreme da curva de potência.

A arquitetura típica de previsão de curto prazo tem duas etapas:
\begin{enumerate}
  \item \textbf{Previsão meteorológica}: modelos NWP globais ou regionais
    fornecem vento em várias alturas, direção, temperatura e pressão.
    Modelos meteorológicos baseados em IA (como GraphCast, GenCast e Aurora)
    passaram a competir com os NWP físicos em várias métricas, a um custo
    computacional muito menor.
  \item \textbf{Modelo de conversão}: um modelo de ML (em geral gradient
    boosting ou redes neurais) aprende a relação entre a meteorologia
    prevista e a potência medida do parque. Esse modelo corrige vieses do
    NWP, efeitos de esteira, indisponibilidades e perdas.
\end{enumerate}

O ONS desenvolveu um modelo próprio de previsão eólica voltado ao Nordeste e
ao Sul para reduzir a necessidade de reserva de potência~\cite{onseolico}.

\begin{nota}
\textbf{Cuidado com a limitação de geração (curtailment).} Quando o operador
restringe a geração de um parque (no Brasil, o chamado
\textit{constrained-off}, que se tornou muito relevante para eólicas e solares
no Nordeste), a potência medida deixa de refletir o recurso disponível. Treinar
com esses dados sem filtrá-los ensina o modelo a subestimar a geração. É
preciso marcar e excluir os períodos de restrição, ou modelar a potência
\emph{disponível} em vez da potência \emph{entregue}.
\end{nota}

\subsection{Previsão de geração solar}

A irradiância que chega à superfície é dominada por um componente
determinístico (a posição do sol) e um componente estocástico (nuvens e
aerossóis). A normalização padrão é o \textbf{índice de céu claro}:
\begin{equation}
  k_t = \frac{\mathrm{GHI}_t}{\mathrm{GHI}^{\mathrm{cs}}_t},
\end{equation}
em que $\mathrm{GHI}^{\mathrm{cs}}$ é a irradiância global horizontal
calculada por um modelo de céu claro. Prever $k_t$ em vez de GHI remove a
sazonalidade astronômica e deixa o modelo concentrado no que é difícil: as
nuvens.

As fontes de dados e modelos variam com o horizonte:
\begin{itemize}
  \item \textbf{minutos}: câmeras de céu (\textit{all-sky imagers}) e CNNs que
    estimam o movimento das nuvens;
  \item \textbf{até 6 horas}: imagens de satélite geoestacionário (no Brasil,
    GOES-East) e modelos de vídeo/espaço-temporais (ConvLSTM, U-Net);
  \item \textbf{dia seguinte e além}: NWP como entrada de modelos de ML.
\end{itemize}

Na conversão de irradiância para potência entram temperatura dos módulos
(eficiência cai com o calor), orientação, sombreamento, sujidade e limite de
potência do inversor (\textit{clipping}). Estudos brasileiros recentes usam
LSTM com sucesso para previsão horária de geração fotovoltaica~\cite{pvlstm}.

\subsection{Previsão de preço}

Em mercados de oferta (Europa, EUA), redes profundas estão entre os melhores
métodos para prever preços do mercado do dia seguinte. O Brasil é diferente: o
Preço de Liquidação das Diferenças (PLD) não resulta de ofertas, mas é
calculado pela CCEE a partir da cadeia de modelos de otimização
hidrotérmica (NEWAVE, DECOMP e DESSEM, desenvolvidos pelo CEPEL). Desde 2021 o
PLD é horário e calculado pelo DESSEM.

Isso abre duas estratégias de ML:
\begin{enumerate}
  \item \textbf{Previsão direta do PLD} a partir de variáveis como afluências,
    armazenamento, carga e geração renovável prevista;
  \item \textbf{Emulação dos modelos oficiais} (\textit{surrogate modeling}):
    treinar uma rede para aproximar a saída do DESSEM/DECOMP em função das
    entradas, permitindo rodar milhares de cenários em segundos. É útil para
    análise de risco de comercializadoras.
\end{enumerate}

\subsection{Famílias de modelos}

\begin{table}[h]
\centering
\small
\begin{tabularx}{\textwidth}{l L L}
\toprule
\textbf{Família} & \textbf{Pontos fortes} & \textbf{Limitações} \\
\midrule
Estatísticos (ARIMA, ETS, regressão) & interpretáveis, baratos, bons
  \textit{baselines} & capturam mal não linearidades e muitas covariáveis \\
Gradient boosting (XGBoost, LightGBM) & excelente em dados tabulares com
  covariáveis ricas; padrão da indústria & exige engenharia de atributos
  (defasagens, janelas) \\
RNN (LSTM, GRU) & modelam dependência temporal diretamente & treinamento
  mais lento; nem sempre superam GBM \\
CNN temporais (TCN) & paralelizáveis, campo receptivo longo & menos usadas
  com covariáveis futuras \\
N-BEATS, N-HiTS & MLP profundas com decomposição interpretável & univariadas
  na forma original \\
Transformers (TFT, PatchTST) & covariáveis heterogêneas, múltiplos
  horizontes, quantis (TFT)~\cite{tft} & exigem mais dados e ajuste \\
Foundation models (Chronos, TimesFM, Moirai) & previsão \textit{zero-shot}
  sem treino local~\cite{chronos,timesfm} & suporte limitado a covariáveis;
  custo de inferência \\
\bottomrule
\end{tabularx}
\caption{Famílias de modelos para previsão em energia.}
\end{table}

\paragraph{Foundation models de séries temporais.} São modelos pré-treinados
em grandes coleções de séries temporais heterogêneas, capazes de prever
séries nunca vistas sem retreino. A linha começou com TimeGPT (2023) e seguiu
com Lag-Llama, Chronos, TimesFM e, mais recentemente, Sundial e
Kairos~\cite{fmreview}. Resultados de 2025--2026 em energia:
\begin{itemize}
  \item em benchmarks de carga em vários níveis da rede, Transformers de
    séries temporais superaram os métodos estabelecidos~\cite{loadbenchmark};
  \item foundation models rodam em hardware de consumo com bom desempenho
    zero-shot, mas com sensibilidade ao tamanho do contexto e à mudança de
    distribuição~\cite{tsfmconsumer};
  \item a incorporação de covariáveis (clima, calendário) melhora
    significativamente os modelos e permite explicar as previsões~\cite{tsfmxai};
  \item na baixa tensão, onde as séries são ruidosas e curtas, os modelos são
    promissores para previsão probabilística de picos~\cite{lvpeak};
  \item modelos específicos de energia, como o GridFM, preveem de forma
    conjunta carga, preço, emissões e renováveis, incorporando restrições
    físicas~\cite{gridfm}.
\end{itemize}
O principal benefício prático é o \textbf{arranque a frio}: novos parques,
novos alimentadores ou novos clientes ganham previsões razoáveis no primeiro
dia, antes de haver histórico para treinar um modelo dedicado.

\subsection{Avaliação}

\paragraph{Métricas.}
\begin{itemize}
  \item \textbf{MAE, RMSE}: na unidade da variável; o RMSE penaliza mais os
    erros grandes, relevantes para rampas.
  \item \textbf{MAPE}: popular em carga, mas \emph{inadequado para solar}, que
    vale zero à noite (divisão por zero) e tem valores pequenos no começo e
    no fim do dia. Para renováveis usa-se o erro normalizado pela capacidade
    instalada (nMAE, nRMSE).
  \item \textbf{Skill score}: ganho relativo sobre uma referência ingênua
    (persistência ou climatologia),
    $\mathrm{SS} = 1 - \mathrm{RMSE}_{\text{modelo}}/\mathrm{RMSE}_{\text{ref}}$.
    Sem essa comparação, números absolutos de erro dizem pouco.
  \item \textbf{Probabilísticas}: pinball loss, CRPS, cobertura e largura dos
    intervalos (calibração e nitidez).
  \item \textbf{Métricas de valor}: custo de desvio evitado, reserva poupada,
    receita adicional. É a métrica que interessa ao negócio.
\end{itemize}

\paragraph{Validação.} Dois erros comuns inflam artificialmente os resultados:
\begin{enumerate}
  \item \textbf{Validação cruzada aleatória} em séries temporais vaza o
    futuro para o treino. O correto é \textit{backtesting} com origem móvel
    (\textit{rolling origin}), respeitando a ordem temporal.
  \item \textbf{Vazamento de covariáveis}: usar a meteorologia
    \emph{observada} no lugar da \emph{prevista} disponível no momento da
    decisão. Em produção só existe a previsão, emitida em um horário
    específico. O treino precisa reproduzir exatamente o que estaria
    disponível em cada momento.
\end{enumerate}

\subsection{Desafios}
\begin{itemize}
  \item \textbf{Eventos extremos}: ondas de calor, tempestades e apagões são
    raros no histórico e justamente os mais importantes.
  \item \textbf{Mudanças estruturais}: crescimento da MMGD, eletrificação,
    mudança de hábitos. Modelos precisam de monitoramento de deriva e
    retreino contínuo.
  \item \textbf{Qualidade dos dados}: medições faltantes, erros de telemetria,
    curtailment não marcado.
  \item \textbf{Confiança do usuário}: o operador precisa entender por que a
    previsão mudou. Explicabilidade (importância de covariáveis, decomposição)
    é requisito, não enfeite.
\end{itemize}

\begin{produto}{Plataforma de previsão para renováveis e carga}

\prodsec{Problema}
Geradores eólicos e solares estão expostos a penalidades por desvio,
curtailment e preço de curto prazo. Comercializadoras precisam de previsões
confiáveis para contratar. Distribuidoras e grandes consumidores têm carga
líquida cada vez mais incerta por causa da MMGD. Muitos desses agentes usam
planilhas, previsões genéricas de fornecedores estrangeiros ou persistência.

\prodsec{Usuários e personas}
\begin{itemize}
  \item \textbf{Analista de operação comercial de um gerador}: precisa da
    previsão de geração por parque para a programação e para declarações ao
    operador.
  \item \textbf{Trader de comercializadora}: precisa de previsões de geração
    agregada, carga e PLD, com incerteza, para decidir posições.
  \item \textbf{Planejador de distribuidora}: precisa da carga por
    subestação/alimentador e da MMGD por região.
\end{itemize}

\prodsec{Proposta de valor}
Previsões probabilísticas por ativo, com horizonte de minutos a 15 dias,
disponíveis via API e painel, calibradas para o contexto brasileiro (NWP e
satélite sobre o território, calendário nacional, marcação de curtailment),
com ganho demonstrável sobre a referência que o cliente já usa.

\prodsec{Escopo do MVP}
\begin{enumerate}
  \item Previsão de geração eólica e solar do dia seguinte e intradiária
    para parques com telemetria disponível.
  \item Quantis (P10, P50, P90) e não só valor pontual.
  \item Painel com comparação diária previsto versus realizado e skill score
    contra persistência.
  \item API REST para integração com sistemas do cliente.
  \item Onboarding de novos parques com foundation model zero-shot, trocado
    por modelo ajustado quando houver histórico.
\end{enumerate}

\prodsec{Evolução (roadmap)}
\begin{itemize}
  \item \textbf{V2}: nowcasting com imagens de satélite (0--6\,h);
    detecção automática de curtailment e indisponibilidade.
  \item \textbf{V3}: previsão de PLD e de carga por submercado; cenários
    conjuntos (geração, carga, preço) para análise de risco.
  \item \textbf{V4}: otimização de ofertas e de despacho de baterias acoplada
    às previsões (da previsão à decisão).
\end{itemize}

\prodsec{Arquitetura de referência}
\begin{itemize}
  \item \textbf{Ingestão}: NWP (múltiplos provedores), satélite, telemetria
    SCADA dos parques, dados abertos do ONS e da CCEE.
  \item \textbf{Feature store} com versão temporal (``o que se sabia em cada
    instante'') para evitar vazamento.
  \item \textbf{Treino} por ativo ou global (um modelo para vários parques),
    com seleção automática por backtesting.
  \item \textbf{Serviço de inferência} agendado conforme a chegada de cada
    rodada de NWP.
  \item \textbf{Monitoramento} de erro em produção, deriva e alertas de dados
    faltantes.
\end{itemize}

\prodsec{Métricas de sucesso}
\begin{itemize}
  \item Técnicas: nMAE e CRPS por ativo; skill score sobre persistência e
    sobre o fornecedor atual do cliente.
  \item De negócio: redução de custo de desvio (R\$/MWh), receita
    adicional, horas de analista economizadas.
  \item De produto: tempo de onboarding de um novo parque, disponibilidade da
    API, retenção de clientes.
\end{itemize}

\prodsec{Modelo de negócio}
Assinatura por MW monitorado ou por ativo, com camadas (dia seguinte;
intradiário; nowcasting; preço). Um piloto gratuito ou de baixo custo em
modo sombra (\textit{shadow mode}), comparando com a previsão atual do cliente,
é o principal instrumento de venda: o ganho aparece em números do próprio
cliente.

\prodsec{Riscos}
\begin{itemize}
  \item Dependência de NWP de terceiros (custo, disponibilidade, mudanças de
    versão que alteram a distribuição de entrada).
  \item Acesso à telemetria do cliente (burocracia, segurança da informação).
  \item Mercado com concorrentes internacionais estabelecidos; a
    diferenciação precisa vir de acurácia local e integração com o contexto
    regulatório brasileiro.
\end{itemize}
\end{produto}

% ===========================================================================
\section{Operação e otimização da rede}
% ===========================================================================

\subsection{Contexto}

A operação de um sistema elétrico em tempo real envolve resolver
repetidamente problemas de análise e otimização:
\begin{itemize}
  \item \textbf{Fluxo de potência}: dado o estado de geração e carga,
    calcular tensões e fluxos em toda a rede;
  \item \textbf{Análise de contingências}: verificar se o sistema continua
    seguro após a perda de qualquer elemento (critério $N-1$), o que exige
    milhares de fluxos de potência;
  \item \textbf{Fluxo de potência ótimo (OPF)}: escolher o despacho de menor
    custo que respeite todos os limites físicos;
  \item \textbf{Estimação de estado}: reconstruir o estado mais provável da
    rede a partir de medições ruidosas e incompletas;
  \item \textbf{Controle}: tensão e reativos, redespacho, chaveamento
    (topologia), baterias.
\end{itemize}

Com mais renováveis, esses problemas precisam ser resolvidos com mais
frequência, para mais cenários e com mais incerteza. Os solvers clássicos
são confiáveis, mas o custo computacional cresce rapidamente com o
tamanho da rede e o número de cenários. ML entra aqui principalmente como
\textbf{acelerador} e como \textbf{apoio à decisão}, não como substituto dos
métodos físicos.

\subsection{Formulação: fluxo de potência AC}

Para cada barra $i$ da rede, as injeções de potência ativa e reativa são
\begin{align}
  P_i &= \sum_{j} |V_i||V_j| \big(G_{ij}\cos\theta_{ij} + B_{ij}\sin\theta_{ij}\big), \\
  Q_i &= \sum_{j} |V_i||V_j| \big(G_{ij}\sin\theta_{ij} - B_{ij}\cos\theta_{ij}\big),
\end{align}
em que $V_i$ é o fasor de tensão, $\theta_{ij} = \theta_i - \theta_j$ e
$G_{ij} + jB_{ij}$ são os elementos da matriz de admitância nodal. O sistema é
não linear e resolvido por Newton--Raphson.

O OPF adiciona uma função objetivo e restrições:
\begin{equation}
\begin{aligned}
  \min_{P^g, Q^g, V, \theta} \quad & \sum_{g} c_g(P^g_g) \\
  \text{s.a.} \quad & \text{equações de fluxo de potência}, \\
  & V_i^{\min} \le |V_i| \le V_i^{\max}, \\
  & P^{g,\min}_g \le P^g_g \le P^{g,\max}_g, \\
  & |S_{ij}| \le S_{ij}^{\max}.
\end{aligned}
\end{equation}
É um problema não convexo e, em sistemas grandes com restrições de
segurança, caro de resolver dentro dos prazos de tempo real.

\subsection{Redes neurais em grafos (GNNs)}

Uma rede elétrica é naturalmente um grafo: barras são nós, linhas e
transformadores são arestas. GNNs processam esse tipo de dado por
\textit{message passing}:
\begin{equation}
  \mathbf{h}_i^{(k+1)} = \phi\Big(\mathbf{h}_i^{(k)},\;
  \bigoplus_{j \in \mathcal{N}(i)} \psi\big(\mathbf{h}_i^{(k)}, \mathbf{h}_j^{(k)}, \mathbf{e}_{ij}\big)\Big),
\end{equation}
em que $\mathbf{h}_i$ é a representação da barra $i$ (carga, geração, tipo),
$\mathbf{e}_{ij}$ são atributos da linha (impedância, limite) e $\bigoplus$ é
uma agregação invariante a permutações (soma, média). Após $K$ camadas, cada
barra ``enxerga'' sua vizinhança de $K$ saltos.

Essa estrutura traz duas vantagens em relação a MLPs que recebem um vetor
fixo:
\begin{enumerate}
  \item \textbf{Generalização topológica}: a mesma GNN funciona quando uma
    linha sai de operação ou quando a rede muda, porque os parâmetros são
    compartilhados entre nós e arestas.
  \item \textbf{Escalabilidade}: o custo cresce com o número de arestas, e
    modelos treinados em redes menores podem ser aplicados a redes maiores.
\end{enumerate}

\paragraph{Aplicações principais.}
\begin{itemize}
  \item \textbf{Surrogate de fluxo de potência}: prever tensões e fluxos
    diretamente, em milissegundos, para triagem rápida de contingências.
    Resultados recentes mostram GNNs com boa acurácia e escalando para
    sistemas maiores com tempo de computação muito menor que o dos solvers
    tradicionais~\cite{gnnpf}.
  \item \textbf{Warm start de OPF}: a GNN fornece um ponto inicial próximo da
    solução, e o solver clássico converge em menos iterações, preservando a
    garantia de factibilidade~\cite{gnnopfinit}.
  \item \textbf{Reconfiguração de rede}: prever a melhor topologia antes da
    etapa de otimização, reduzindo o espaço de busca
    combinatório~\cite{gnnreconf}.
\end{itemize}

\subsection{Aprender a otimizar e o problema da factibilidade}

Treinar uma rede para mapear diretamente entradas (carga, custos) em
soluções de OPF (\textit{optimization proxies}) esbarra em um problema: a rede
não garante que a solução respeite as restrições. Uma violação de limite
térmico ou de tensão não é um ``pequeno erro''; é uma solução inaceitável. As
principais estratégias são:
\begin{itemize}
  \item \textbf{Penalização} das violações na função de perda (simples, sem
    garantia);
  \item \textbf{Dualidade Lagrangiana} durante o treino, ajustando pesos das
    restrições;
  \item \textbf{Correção e completamento}: a rede prevê parte das variáveis e
    as demais são obtidas resolvendo as igualdades, seguido de passos de
    correção das desigualdades (como no DC3~\cite{dc3});
  \item \textbf{Projeção} da saída no conjunto factível, ou uso da previsão
    apenas como warm start de um solver exato.
\end{itemize}
Na prática industrial, a abordagem mais aceita é a última: o ML acelera, o
solver certifica.

\subsection{Aprendizado por reforço (RL)}

Problemas de controle sequencial, como decidir chaveamentos, redespachos ou
ciclos de baterias ao longo do tempo, podem ser formulados como processos de
decisão de Markov. O agente observa o estado $s_t$ (fluxos, tensões,
previsões), escolhe uma ação $a_t$ e recebe uma recompensa $r_t$ (custo
evitado, margem de segurança). O objetivo é
\begin{equation}
  \max_\pi \; \mathbb{E}_\pi\Big[\sum_{t=0}^{T} \gamma^t r(s_t, a_t)\Big].
\end{equation}

\paragraph{Controle de topologia.} A competição \textit{Learning to Run a
Power Network} (L2RPN), organizada pela RTE (operador francês), popularizou o
uso de RL para manter a rede operando via mudanças de topologia em
subestações, redespacho e curtailment~\cite{l2rpn}. Agentes que combinam GNNs
com RL (\textit{graph RL}), às vezes pré-treinados por aprendizado por
imitação e depois refinados com PPO, lidam melhor com mudanças de topologia
causadas por clima extremo ou manutenção~\cite{graphrlsurvey,ppognn}.

\paragraph{Baterias e microrredes.} RL é usado para decidir quando carregar e
descarregar baterias considerando preço, previsão de geração e degradação.
Surrogates baseados em PINNs permitem treinar agentes sem interagir com o
sistema real~\cite{rlpinn}.

\paragraph{RL seguro.} Em sistemas físicos, explorar ações ruins tem custo
real. A formulação padrão é o MDP com restrições:
\begin{equation}
  \max_\pi \; \mathbb{E}_\pi\Big[\sum_t \gamma^t r_t\Big]
  \quad \text{s.a.} \quad
  \mathbb{E}_\pi\Big[\sum_t \gamma^t c_t\Big] \le d,
\end{equation}
em que $c_t$ mede violações (sobrecarga, subtensão). Outras técnicas incluem
\textit{shielding} (um módulo verifica e bloqueia ações perigosas antes da
execução) e treino offline a partir de dados históricos de operadores~\cite{saferl}.

\subsection{Redes neurais informadas pela física (PINNs)}

PINNs incorporam equações físicas na função de perda:
\begin{equation}
  \mathcal{L} = \mathcal{L}_{\text{dados}} + \lambda\, \mathcal{L}_{\text{física}},
\end{equation}
em que $\mathcal{L}_{\text{física}}$ penaliza o resíduo das equações (por
exemplo, das equações de fluxo de potência ou de equações diferenciais de
dinâmica de máquinas)~\cite{pinn}. Os benefícios são menor necessidade de dados
rotulados, maior respeito às leis físicas e melhor extrapolação. São usadas em
estimação de estado, análise de estabilidade transitória e como simuladores
para treinar agentes de RL.

\subsection{Desafios}
\begin{itemize}
  \item \textbf{Certificação e responsabilidade}: operadores e reguladores
    exigem garantias; um modelo que acerta 99,9\% dos casos ainda pode falhar
    exatamente na contingência crítica.
  \item \textbf{Generalização}: a rede real muda (novas linhas, novas
    usinas), e modelos treinados em cenários passados podem não cobrir
    situações futuras~\cite{gridgeneralization}.
  \item \textbf{Integração com EMS legados}: sistemas de gerenciamento de
    energia são críticos, fechados e com ciclos de atualização longos.
  \item \textbf{Dados}: modelos detalhados da rede e dados operacionais são
    sensíveis; a pesquisa depende muito de sistemas de teste sintéticos (IEEE,
    PGLib).
  \item \textbf{Fator humano}: a autonomia plena não é realista no curto
    prazo. O papel viável é o de copiloto do operador.
\end{itemize}

\begin{produto}{Copiloto de operação da rede}

\prodsec{Problema}
Operadores de centros de controle (de transmissão ou de distribuição)
precisam tomar decisões rápidas diante de um número crescente de cenários:
variabilidade renovável, curtailment, contingências, restrições de
transmissão. As ferramentas de análise são lentas para explorar muitas
alternativas em tempo real, e a experiência dos operadores mais antigos não
está codificada.

\prodsec{Usuários e personas}
\begin{itemize}
  \item \textbf{Operador de tempo real} de uma transmissora ou do centro de
    operação da distribuição (COD);
  \item \textbf{Engenheiro de estudos elétricos} que prepara a operação do dia
    seguinte;
  \item \textbf{Operador de parques renováveis e baterias} que precisa decidir
    despacho e armazenamento.
\end{itemize}

\prodsec{Proposta de valor}
Triagem de contingências e recomendação de ações corretivas em segundos, com
cada recomendação verificada por um solver físico antes de ser apresentada.
O operador continua no comando; o sistema amplia quantos cenários ele
consegue avaliar.

\prodsec{Escopo do MVP}
\begin{enumerate}
  \item Surrogate de fluxo de potência baseado em GNN, treinado com
    simulações do modelo da rede do cliente.
  \item Triagem rápida de contingências $N-1$ para os próximos intervalos,
    usando previsões de carga e renováveis, com os casos críticos
    confirmados por fluxo de potência completo.
  \item Painel de ``pontos quentes'': quais linhas e barras correm risco e
    quando.
  \item Operação em \textbf{modo sombra}: o sistema roda em paralelo, sem atuar,
    e se registram as recomendações versus as decisões reais.
\end{enumerate}

\prodsec{Evolução (roadmap)}
\begin{itemize}
  \item \textbf{V2}: recomendação de ações corretivas (redespacho,
    chaveamento), ranqueadas e verificadas pelo solver.
  \item \textbf{V3}: agente de RL treinado com decisões históricas e em
    simulação, sugerindo sequências de ações com restrições de segurança.
  \item \textbf{V4}: otimização coordenada de baterias e curtailment em
    tempo real para parques híbridos.
\end{itemize}

\prodsec{Arquitetura de referência}
\begin{itemize}
  \item \textbf{Gêmeo digital} da rede (modelo elétrico sincronizado com a
    topologia atual).
  \item \textbf{Gerador de cenários} para criar dados de treino sintéticos a
    partir do solver físico.
  \item \textbf{Modelos}: GNN de fluxo de potência; ranqueador de
    contingências; política de RL com \textit{shield}.
  \item \textbf{Camada de verificação}: toda saída passa pelo solver antes de
    chegar ao operador.
  \item \textbf{Integração} com SCADA/EMS via protocolos padrão (como ICCP e
    CIM), em modo apenas leitura no início.
\end{itemize}

\prodsec{Métricas de sucesso}
\begin{itemize}
  \item Técnicas: erro de tensão e fluxo do surrogate; \textit{recall} de
    contingências críticas (perder uma crítica é o erro inaceitável);
    aceleração em relação ao solver.
  \item Operacionais: redução do curtailment evitável, redução de violações,
    tempo de resposta a eventos.
  \item De adoção: percentual de recomendações aceitas pelos operadores.
\end{itemize}

\prodsec{Modelo de negócio}
Licença de software por centro de controle, com serviço de implantação e
calibração do gêmeo digital. Para transmissoras e distribuidoras, projetos de
P\&D regulados pela ANEEL são uma porta de entrada natural para pilotos.

\prodsec{Riscos}
\begin{itemize}
  \item Ciclo de venda longo e aversão a risco dos clientes.
  \item Acesso ao modelo da rede e aos dados operacionais.
  \item Requisitos de cibersegurança para qualquer integração com sistemas de
    controle.
  \item Responsabilidade em caso de recomendação errada; o desenho de
    ``humano no controle'' e a verificação física mitigam, mas não eliminam.
\end{itemize}
\end{produto}

% ===========================================================================
\section{Detecção de falhas, proteção e manutenção preditiva}
% ===========================================================================

\subsection{Contexto}

Falhas em equipamentos causam interrupções de fornecimento, danos a
ativos caros e riscos à segurança. No Brasil, a ANEEL mede a continuidade do
fornecimento por indicadores coletivos (DEC, duração, e FEC, frequência
equivalente de interrupção por consumidor) e individuais (DIC, FIC, DMIC), e
a violação dos limites individuais gera compensações pagas aos consumidores.
Reduzir a frequência e a duração das falhas tem, portanto, valor financeiro
direto.

A manutenção evolui em quatro estágios:
\begin{enumerate}
  \item \textbf{Corretiva}: consertar depois de quebrar;
  \item \textbf{Preventiva}: intervenções periódicas por calendário;
  \item \textbf{Baseada em condição}: intervir quando indicadores medidos
    ultrapassam limites;
  \item \textbf{Preditiva}: estimar quando a falha deve ocorrer e agir antes.
\end{enumerate}
ML é o que viabiliza a passagem do terceiro para o quarto estágio em larga
escala.

\subsection{Detecção, classificação e localização de faltas}

Uma falta (curto-circuito) em uma linha produz transitórios de tensão e
corrente registrados por relés digitais e oscilógrafos. A análise envolve três
tarefas~\cite{faultsurvey}:
\begin{itemize}
  \item \textbf{Detecção}: houve falta?
  \item \textbf{Classificação}: que tipo (fase-terra, fase-fase, trifásica
    etc.)?
  \item \textbf{Localização}: a que distância do terminal?
\end{itemize}

\paragraph{Abordagens.}
\begin{itemize}
  \item \textbf{Extração de atributos + classificador}: transformada wavelet
    ou de Fourier sobre os sinais, seguida de SVM, árvores ou redes rasas.
  \item \textbf{Aprendizado de ponta a ponta}: CNNs 1D sobre as formas de onda
    brutas ou CNNs 2D sobre representações tempo-frequência (espectrogramas,
    escalogramas).
  \item \textbf{Dados de PMU}: medições fasoriais sincronizadas em alta taxa
    permitem detectar oscilações e eventos em toda a rede; modelos
    espaço-temporais (GNN + redes temporais) exploram a topologia.
  \item \textbf{Redes de distribuição}: a localização é mais difícil pela
    topologia radial ramificada e pela menor instrumentação; ML combina dados
    de medidores inteligentes, chamadas de clientes e religadores.
\end{itemize}

Um detalhe importante é que dados reais de faltas são escassos. Por isso,
os modelos são frequentemente treinados com \textbf{simulações de transitórios
eletromagnéticos} (ATP/EMTP, PSCAD), o que cria a lacuna entre simulação e
realidade (\textit{sim-to-real}), um dos principais riscos dessas aplicações.

\subsection{Manutenção preditiva de ativos}

\paragraph{Transformadores de potência.} São dos ativos mais caros e críticos.
A \textbf{análise de gases dissolvidos} (DGA) no óleo isolante é a principal
fonte de diagnóstico: diferentes tipos de falha interna (descargas parciais,
arcos, sobreaquecimento) geram proporções características de hidrogênio,
metano, etileno, acetileno e outros gases. Métodos clássicos, como o triângulo
de Duval, são regras gráficas; ML permite combinar DGA com histórico de carga,
temperatura, idade e ensaios para produzir um \textit{índice de saúde} e uma
probabilidade de falha.

\paragraph{Turbinas eólicas.} Os sistemas SCADA das turbinas registram, a cada
10 minutos, dezenas de variáveis (temperaturas de rolamentos e da caixa
multiplicadora, potência, vento, rotação). A técnica mais difundida é o
\textbf{modelo de comportamento normal}:
\begin{enumerate}
  \item treina-se um modelo que prevê uma variável (por exemplo, temperatura
    do rolamento) a partir das demais, usando apenas períodos saudáveis;
  \item em operação, o resíduo $r_t = y_t - \hat{y}_t$ é monitorado;
  \item um desvio persistente do resíduo indica degradação, muitas vezes
    semanas antes da falha.
\end{enumerate}
Sensores de vibração em alta frequência permitem diagnóstico mais fino
(rolamentos, engrenagens) com CNNs sobre espectros.

\paragraph{Usinas fotovoltaicas.} Detecção de strings com baixo desempenho
comparando com vizinhas e com a irradiância esperada; falhas de inversores;
degradação e sujidade (que também orienta a programação de limpeza).

\paragraph{Vida útil remanescente (RUL).} Quando há falhas suficientes no
histórico, pode-se estimar o tempo até a falha com modelos de sobrevivência
(Cox, \textit{random survival forests}, redes de sobrevivência) ou regressão
direta com redes recorrentes. O resultado alimenta a otimização do plano de
manutenção, que pondera custo da intervenção, probabilidade de falha e
consequência.

\subsection{Visão computacional para inspeção}

\begin{itemize}
  \item \textbf{Linhas de transmissão e distribuição}: detecção de isoladores
    danificados, amortecedores faltantes, corrosão e ninhos em imagens de
    drones e helicópteros, com detectores de objetos (família YOLO, DETR).
  \item \textbf{Vegetação}: nuvens de pontos LiDAR e imagens de satélite para
    medir distância entre copas e condutores e priorizar podas, uma das
    principais causas de falhas na distribuição.
  \item \textbf{Termografia}: pontos quentes em conexões, chaves e painéis
    solares.
  \item \textbf{Eletroluminescência}: microfissuras em células fotovoltaicas.
\end{itemize}
O ganho aqui é de escala: uma inspeção que exigiria equipes percorrendo
milhares de quilômetros passa a ser um problema de triagem de imagens, com
humanos revisando apenas os casos sinalizados.

\subsection{Desafios de dados}
\begin{itemize}
  \item \textbf{Desbalanceamento extremo}: falhas são raras. Métricas como
    acurácia não servem; usam-se precisão, recall, curva PR, e custo de falsos
    positivos (visita desnecessária) versus falsos negativos (falha não
    prevista).
  \item \textbf{Rótulos pobres}: ordens de serviço em texto livre, causas
    registradas de forma inconsistente. Processamento de linguagem natural
    (hoje com LLMs) ajuda a estruturar o histórico de manutenção.
  \item \textbf{Dados dispersos}: SCADA, sistemas de manutenção, laboratório de
    óleo e inspeções em sistemas diferentes, sem chave comum.
  \item \textbf{Heterogeneidade dos ativos}: fabricantes, modelos e idades
    diferentes reduzem o volume de dados comparáveis. Aprendizado por
    transferência e modelos semi-supervisionados ajudam.
\end{itemize}

\begin{produto}{Plataforma de saúde de ativos (APM)}

\prodsec{Problema}
Transmissoras, distribuidoras e geradores mantêm grandes frotas de ativos com
planos de manutenção baseados principalmente em calendário. Isso gera
manutenção desnecessária em ativos saudáveis e falhas inesperadas em ativos
degradados, com custo de reparo, perda de receita, compensações por
continuidade e risco à segurança.

\prodsec{Usuários e personas}
\begin{itemize}
  \item \textbf{Engenheiro de manutenção}: quer saber quais ativos precisam de
    atenção e por quê;
  \item \textbf{Gestor de ativos}: quer priorizar orçamento de manutenção e
    substituição com base em risco;
  \item \textbf{Equipe de inspeção}: quer rotas e ordens de serviço
    priorizadas.
\end{itemize}

\prodsec{Proposta de valor}
Um índice de saúde e de risco por ativo, atualizado continuamente, com a
explicação dos sinais que o motivaram, integrado ao sistema de ordens de
serviço. A decisão passa de ``manter tudo a cada $N$ meses'' para ``manter o
que está degradando, quando está degradando''.

\prodsec{Escopo do MVP}
Começar com \textbf{uma classe de ativo} em que dados e valor estejam claros:
\begin{itemize}
  \item \textbf{Opção A (geração eólica)}: modelos de comportamento normal
    sobre SCADA de turbinas para caixa multiplicadora, gerador e rolamentos,
    com alertas e painel por parque;
  \item \textbf{Opção B (transmissão/distribuição)}: índice de saúde de
    transformadores a partir de DGA, carregamento, idade e histórico de
    ensaios.
\end{itemize}
Em ambos os casos: ranking de ativos por risco, explicação dos alertas e
registro do \textit{feedback} do engenheiro (alerta útil ou falso), que vira
rótulo para o próximo ciclo de treino.

\prodsec{Evolução (roadmap)}
\begin{itemize}
  \item \textbf{V2}: módulo de visão computacional para inspeção por drone
    (linhas, subestações, painéis solares).
  \item \textbf{V3}: estimativa de vida útil remanescente e otimização do
    plano de manutenção sob restrição de orçamento e equipes.
  \item \textbf{V4}: análise automática de oscilografias para
    classificação e localização de faltas, integrada à recomposição.
\end{itemize}

\prodsec{Arquitetura de referência}
\begin{itemize}
  \item \textbf{Ingestão} de SCADA/historiadores (como PI), laudos de
    laboratório, sistemas de gestão de manutenção e imagens.
  \item \textbf{Cadastro unificado de ativos} (a chave que liga os dados, na
    prática o maior trabalho do projeto).
  \item \textbf{Modelos} por classe de ativo, com camada comum de alertas e
    explicabilidade.
  \item \textbf{Fluxo de trabalho}: alerta $\rightarrow$ análise do engenheiro
    $\rightarrow$ ordem de serviço $\rightarrow$ resultado da inspeção
    $\rightarrow$ rótulo.
\end{itemize}

\prodsec{Métricas de sucesso}
\begin{itemize}
  \item Técnicas: precisão dos alertas (proporção de alertas confirmados),
    antecedência média entre alerta e falha, recall de falhas.
  \item Operacionais: redução de falhas não programadas, redução de
    manutenções desnecessárias, melhora de DEC/FEC e redução de compensações.
  \item Financeiras: custo de manutenção por ativo e disponibilidade da
    geração.
\end{itemize}

\prodsec{Modelo de negócio}
Assinatura por ativo monitorado (por turbina, por transformador) mais serviço
de implantação. Contratos por resultado (parte do valor atrelada à redução de
falhas) são possíveis quando há linha de base clara.

\prodsec{Riscos}
\begin{itemize}
  \item \textbf{Fadiga de alertas}: muitos falsos positivos fazem os usuários
    ignorarem o sistema. Poucos alertas bons valem mais que muitos alertas
    medíocres.
  \item Qualidade e integração dos dados legados.
  \item Concorrência de fabricantes (OEMs), que oferecem monitoramento dos
    próprios equipamentos; a diferenciação é ser multifabricante.
\end{itemize}
\end{produto}

% ===========================================================================
\section{Distribuição e consumidor final}
% ===========================================================================

\subsection{Contexto}

A distribuição é o segmento em que a transformação é mais visível: a
MMGD cresceu muito após o marco legal da geração distribuída (Lei
14.300/2022), a medição inteligente avança, e novos serviços (tarifas
horárias, resposta da demanda, veículos elétricos, mercado livre para
consumidores menores) mudam a relação com o consumidor. Ao mesmo tempo,
problemas antigos, como as perdas não técnicas, continuam com grande peso
financeiro para as distribuidoras brasileiras.

\subsection{Perdas não técnicas}

\paragraph{O problema.} As perdas de uma distribuidora são a diferença entre a
energia injetada na rede e a energia faturada. Parte são \emph{perdas
técnicas} (efeito Joule em condutores e transformadores); o restante são
\emph{perdas não técnicas}: furto (ligações clandestinas), fraude (adulteração
de medidores), erros de medição e de cadastro. A ANEEL define, nas revisões
tarifárias, níveis regulatórios de perdas que podem ser repassados à tarifa;
o que excede esse nível é prejuízo da distribuidora. Em algumas áreas de
concessão, as perdas não técnicas são uma fração muito relevante do mercado
de baixa tensão.

\paragraph{Formulação como classificação.} Para cada unidade consumidora $u$,
estima-se a probabilidade de irregularidade
$p_u = P(\text{irregular} \mid \mathbf{x}_u)$ a partir de atributos como:
\begin{itemize}
  \item \textbf{perfil de consumo}: quedas abruptas, consumo zero, baixa
    variância, desvio em relação a vizinhos de mesmo perfil;
  \item \textbf{histórico}: inspeções anteriores, irregularidades
    confirmadas, trocas de medidor, reclamações;
  \item \textbf{cadastro}: classe, atividade, carga declarada, tipo de
    medidor;
  \item \textbf{contexto espacial}: taxa de irregularidade na vizinhança e no
    transformador;
  \item \textbf{balanço energético}: quando há medição no transformador,
    \begin{equation}
      E^{\text{PNT}}_{k} = E^{\text{trafo}}_{k} - \sum_{u \in k} E^{\text{medida}}_{u}
      - \hat{E}^{\text{técnica}}_{k},
    \end{equation}
    indica quanto se perde sob cada transformador $k$ e ajuda a localizar
    onde procurar.
\end{itemize}
A saída é uma lista priorizada de alvos para inspeção.

\paragraph{Viés de seleção.} Os rótulos vêm de inspeções, e as inspeções não
são aleatórias: historicamente, foram feitas onde havia suspeita. O modelo
aprende, portanto, a reproduzir os critérios antigos de seleção, e não
necessariamente a identificar fraude~\cite{ntlreview}. Mitigações:
\begin{itemize}
  \item reservar parte das inspeções para amostragem aleatória ou exploratória
    (equilíbrio entre explorar e explorar o conhecido);
  \item aprendizado com exemplos positivos e não rotulados (\textit{PU
    learning}), tratando não inspecionados como não rotulados e não como
    regulares;
  \item modelos de detecção de anomalias não supervisionados como
    complemento.
\end{itemize}

\paragraph{Métricas.} Como a capacidade de inspeção é limitada, a métrica
relevante é a \textbf{taxa de acerto no topo da lista} (precisão em $k$) e a
\textbf{energia recuperada por inspeção}, e não métricas globais como AUC.

\subsection{Desagregação de carga (NILM)}

\textit{Non-Intrusive Load Monitoring} estima o consumo de cada aparelho a
partir de uma única medição agregada~\cite{hart}:
\begin{equation}
  P(t) = \sum_{i=1}^{N} p_i(t) + \varepsilon(t).
\end{equation}
É um problema de separação de fontes mal posto: muitas combinações de
aparelhos explicam o mesmo agregado.

\paragraph{Métodos.}
\begin{itemize}
  \item \textbf{Clássicos}: detecção de eventos (degraus de potência) e
    modelos ocultos de Markov fatoriais.
  \item \textbf{Aprendizado profundo}: redes que mapeiam uma janela do
    agregado no consumo de um aparelho, seja a janela inteira
    (\textit{seq2seq}) ou o ponto central (\textit{seq2point})~\cite{kelly,seq2point}.
    Treinadas em bases públicas como REDD e UK-DALE.
\end{itemize}

\paragraph{Limitação prática.} A qualidade depende muito da taxa de
amostragem. Com medidores que registram a cada 15 minutos ou 1 hora, só é
possível identificar grandes cargas (chuveiro, ar-condicionado, aquecedor,
carregador de veículo elétrico). Para desagregação fina são necessários
dados em segundos ou mais rápidos, obtidos por dispositivos adicionais.

\subsection{Rede de distribuição com MMGD e veículos elétricos}

\begin{itemize}
  \item \textbf{Previsão na baixa tensão}: séries por transformador ou por
    consumidor são ruidosas e pouco previsíveis individualmente. Previsão
    probabilística de picos, inclusive com foundation models, é o caminho
    mais promissor~\cite{lvpeak}.
  \item \textbf{Capacidade de hospedagem}: quanto de MMGD ou de carregadores um
    alimentador suporta sem violar tensão e carregamento. Surrogates de ML
    aceleram estudos que exigiriam milhares de simulações.
  \item \textbf{Identificação de topologia e de fase}: cadastros de qual
    consumidor está em qual transformador e fase são frequentemente errados;
    correlação entre séries de tensão de medidores inteligentes permite
    corrigi-los.
  \item \textbf{Detecção de MMGD não cadastrada} a partir do perfil de
    consumo.
\end{itemize}

\subsection{Resposta da demanda e eficiência}

\begin{itemize}
  \item \textbf{Tarifas horárias}: com a Tarifa Branca e, no futuro, tarifas
    mais granulares, o consumidor precisa entender quando deslocar consumo.
    Modelos de previsão e recomendação personalizam essa orientação.
  \item \textbf{Controle de cargas flexíveis}: ar-condicionado, aquecimento de
    água, carregamento de veículos e baterias podem ser programados por
    controle preditivo baseado em modelo (MPC) ou RL para minimizar custo
    respeitando o conforto.
  \item \textbf{Edificações comerciais e indústria}: previsão de consumo e
    detecção de anomalias identificam desperdício (equipamentos ligados fora
    do horário, degradação de sistemas de climatização). O caso da DeepMind,
    com redução de até 40\% da energia de refrigeração de data centers,
    ilustra o potencial~\cite{deepmindcooling}.
  \item \textbf{Segmentação de clientes}: agrupamento de curvas de carga para
    desenho de tarifas, campanhas de eficiência e ofertas no mercado livre.
\end{itemize}

\subsection{Privacidade e regulação}

Dados de consumo em alta resolução revelam hábitos (horários, presença,
aparelhos). No Brasil, são dados pessoais sob a LGPD (Lei 13.709/2018).
Isso exige base legal para o tratamento, minimização, transparência e
segurança. Técnicas como agregação, aprendizado federado e privacidade
diferencial reduzem a exposição. O avanço da medição inteligente no Brasil
também é mais gradual do que na Europa, o que limita a disponibilidade de
dados de alta resolução em larga escala.

\begin{produto}{Inteligência de perdas e de rede para distribuidoras}

\prodsec{Problema}
Distribuidoras perdem receita com perdas não técnicas acima do nível
regulatório e fazem inspeções com baixa taxa de acerto. Ao mesmo tempo,
precisam lidar com uma rede de baixa tensão cada vez mais complexa (MMGD,
veículos elétricos) com cadastros imprecisos e pouca observabilidade.

\prodsec{Usuários e personas}
\begin{itemize}
  \item \textbf{Gerente de recuperação de energia}: quer maximizar energia
    recuperada com o orçamento de inspeções disponível;
  \item \textbf{Equipe de campo}: quer rotas e alvos priorizados, com
    registro fácil do resultado;
  \item \textbf{Engenheiro de planejamento da distribuição}: quer saber onde a
    rede vai saturar com MMGD e cargas novas.
\end{itemize}

\prodsec{Proposta de valor}
Mais energia recuperada por inspeção, com uma lista priorizada e explicada de
alvos que melhora a cada ciclo com o resultado das próprias inspeções. Em
seguida, a mesma base de dados sustenta a visibilidade da rede de baixa
tensão.

\prodsec{Escopo do MVP}
\begin{enumerate}
  \item Modelo de priorização de alvos de inspeção com dados de faturamento,
    cadastro e histórico de inspeções.
  \item Explicação por alvo (por que este cliente foi indicado).
  \item Ciclo fechado: resultado da inspeção retorna como rótulo; parte das
    inspeções é reservada para exploração, corrigindo o viés de seleção.
  \item Painel de resultados: taxa de acerto e energia recuperada comparadas
    com o método atual da distribuidora (teste A/B por região ou por lote).
\end{enumerate}

\prodsec{Evolução (roadmap)}
\begin{itemize}
  \item \textbf{V2}: balanço energético por transformador onde houver
    medição, combinando localização (onde se perde) com priorização (quem
    investigar).
  \item \textbf{V3}: correção de cadastro de topologia e fase, detecção de
    MMGD não cadastrada, previsão de carregamento de transformadores.
  \item \textbf{V4}: estudos de capacidade de hospedagem e produto voltado ao
    consumidor (desagregação, recomendações de tarifa e eficiência).
\end{itemize}

\prodsec{Arquitetura de referência}
\begin{itemize}
  \item \textbf{Ingestão} do sistema comercial (faturamento, cadastro),
    sistema de gestão de ordens de serviço, medição inteligente onde houver, e
    dados georreferenciados da rede.
  \item \textbf{Motor de priorização} com modelos supervisionados (gradient
    boosting é o padrão) e componente de anomalias não supervisionado.
  \item \textbf{Aplicativo de campo} para registrar resultados das inspeções
    de forma estruturada.
  \item \textbf{Governança de dados} conforme a LGPD (controle de acesso,
    anonimização nas análises agregadas, registro de tratamento).
\end{itemize}

\prodsec{Métricas de sucesso}
\begin{itemize}
  \item Técnicas: precisão no topo da lista; calibração das probabilidades.
  \item De negócio: energia recuperada (MWh e R\$) por inspeção; redução do
    índice de perdas não técnicas; custo por inspeção produtiva.
  \item De adoção: percentual das inspeções geradas pelo sistema;
    percentual de resultados registrados.
\end{itemize}

\prodsec{Modelo de negócio}
Assinatura mais componente variável ligado à energia recuperada
(\textit{success fee}). O ganho é mensurável em reais e aparece rápido, o que
facilita a venda. Contratos de P\&D regulados pela ANEEL também podem
financiar os módulos de rede mais inovadores.

\prodsec{Riscos}
\begin{itemize}
  \item Qualidade dos rótulos e viés de seleção; sem ciclo exploratório o
    modelo se torna autorreferente.
  \item Sensibilidade social e jurídica: a indicação de irregularidade precisa
    de procedimento justo; o modelo prioriza inspeções, não constitui prova.
  \item Viés geográfico ou socioeconômico, que precisa ser monitorado para
    evitar concentração injustificada de inspeções em determinados grupos.
  \item Conformidade com a LGPD e segurança dos dados de consumidores.
\end{itemize}
\end{produto}

% ===========================================================================
\section{Síntese comparativa}
% ===========================================================================

\begin{table}[h]
\centering
\small
\begin{tabularx}{\textwidth}{l L L L l}
\toprule
\textbf{Frente} & \textbf{Técnicas centrais} & \textbf{Produto} &
  \textbf{Cliente principal} & \textbf{Maturidade} \\
\midrule
Previsão & GBM, LSTM, Transformers, foundation models &
  Plataforma de previsão & geradores, comercializadoras, distribuidoras &
  alta \\
Operação da rede & GNN, PINN, RL seguro & Copiloto de operação &
  transmissoras, distribuidoras, operadores de renováveis & média/baixa \\
Falhas e manutenção & CNN, modelos de comportamento normal, sobrevivência &
  Plataforma de saúde de ativos & geradores, transmissoras, distribuidoras &
  alta \\
Distribuição e consumidor & GBM, PU learning, NILM, MPC &
  Inteligência de perdas e de rede & distribuidoras & alta (perdas),
  média (rede) \\
\bottomrule
\end{tabularx}
\caption{Comparação entre as quatro frentes.}
\end{table}

\paragraph{Padrões comuns entre as quatro frentes.}
\begin{enumerate}
  \item \textbf{Dados antes de modelos}: em todas as frentes, o maior esforço
    de implantação é integrar e limpar dados de sistemas legados, e não
    treinar redes.
  \item \textbf{Física como aliada}: modelos que respeitam as leis físicas
    (PINNs, verificação por solver, normalizações como o índice de céu claro)
    generalizam melhor e ganham a confiança dos engenheiros.
  \item \textbf{Humano no controle}: os produtos viáveis no curto prazo apoiam
    decisões em vez de tomá-las sozinhos, e aprendem com o
    \textit{feedback} dos usuários.
  \item \textbf{Modo sombra como estratégia de venda}: rodar em paralelo e
    mostrar o ganho com dados do próprio cliente reduz o risco percebido.
  \item \textbf{Ordem de entrada sugerida}: previsão e perdas não técnicas têm
    valor rápido e mensurável e servem de porta de entrada; manutenção
    preditiva vem em seguida; o copiloto de operação é a aposta de maior
    prazo e maior barreira de entrada.
\end{enumerate}

% ===========================================================================
% Referências
% ===========================================================================
\begin{thebibliography}{99}
\small

\bibitem{energyinformatics2025}
Machine learning applications in energy systems: current trends, challenges,
and research directions. \textit{Energy Informatics}, 2025.
\url{https://energyinformatics.springeropen.com/articles/10.1186/s42162-025-00524-6}

\bibitem{trustworthy2024}
Trustworthy artificial intelligence in the energy sector: landscape analysis
and evaluation framework. arXiv:2412.07782, 2024.

\bibitem{deepmindwind}
DeepMind. Machine learning can boost the value of wind energy, 2019. Caso
resumido em
\url{https://bestpractice.ai/ai-use-cases/case-studies/energy-utilities/deepmind-increases-value-of-wind-power-by-20-by-predicting-supply-36-hours-in-advance}

\bibitem{deepmindcooling}
DeepMind. DeepMind AI reduces Google data centre cooling bill by 40\%, 2016.
\url{https://deepmind.google/blog/deepmind-ai-reduces-google-data-centre-cooling-bill-by-40/}

\bibitem{onssolar}
ONS. Avaliação e diagnóstico do sistema de previsão da irradiância solar e
geração. Apresentação em workshop.
\url{https://www.ons.org.br/AcervoDigitalDocumentosEPublicacoes/Apresentacao_Workshop.pdf}

\bibitem{onseolico}
Desenvolvimento e implantação no ONS de um modelo de previsão de geração
eólica. ABEEólica.
\url{https://abeeolica.org.br/wp-content/uploads/2017/07/Art.pdf}

\bibitem{pvlstm}
Previsão horária de geração fotovoltaica utilizando deep learning.
\textit{Observatorio de la Economía Latinoamericana}.
\url{https://ojs.observatoriolatinoamericano.com/ojs/index.php/olel/article/view/11654}

\bibitem{tft}
B.~Lim, S.~Ö.~Arık, N.~Loeff, T.~Pfister. Temporal Fusion Transformers for
interpretable multi-horizon time series forecasting. \textit{International
Journal of Forecasting}, 2021.

\bibitem{chronos}
A.~F.~Ansari et al. Chronos: learning the language of time series.
\textit{Transactions on Machine Learning Research}, 2024.

\bibitem{timesfm}
A.~Das, W.~Kong, R.~Sen, Y.~Zhou. A decoder-only foundation model for
time-series forecasting. \textit{ICML}, 2024.

\bibitem{fmreview}
Foundation models for clean energy forecasting: a comprehensive review.
arXiv:2507.23147, 2025.

\bibitem{loadbenchmark}
A benchmark for electrical load forecasting across grid levels: time-series
transformers outperform established methods. arXiv:2607.15705, 2026.

\bibitem{tsfmconsumer}
Time series foundation models for energy load forecasting on consumer
hardware: a multi-dimensional zero-shot benchmark. arXiv:2602.10848, 2026.

\bibitem{tsfmxai}
Explainable load forecasting with covariate-informed time series foundation
models. arXiv:2604.28149, 2026.

\bibitem{lvpeak}
Probabilistic low-voltage peak load forecasting with time series foundation
models evaluated on application-oriented metrics. arXiv:2607.01966, 2026.

\bibitem{gridfm}
GridFM: a physics-informed foundation model for multi-task energy forecasting
using real-time NYISO data. \textit{Energies}, 19(2):357, 2026.
\url{https://doi.org/10.3390/en19020357}

\bibitem{gnnpf}
Graph neural networks for efficient AC power flow prediction in power grids.
arXiv:2502.05702, 2025.

\bibitem{gnnopfinit}
Initial estimate of AC optimal power flow with graph neural networks.
\textit{Electric Power Systems Research}, 2024.

\bibitem{gnnreconf}
Graph neural network-accelerated network-reconfigured optimal power flow.
arXiv:2410.17460, 2024.

\bibitem{dc3}
P.~L.~Donti, D.~Rolnick, J.~Z.~Kolter. DC3: a learning method for optimization
with hard constraints. \textit{ICLR}, 2021.

\bibitem{l2rpn}
A.~Marot et al. Learning to run a power network challenge for training
topology controllers. \textit{Electric Power Systems Research}, 2020.

\bibitem{graphrlsurvey}
Graph reinforcement learning for power grids: a comprehensive survey.
\textit{Energy and AI}, 2025.

\bibitem{ppognn}
Proximal policy optimization with graph neural networks for optimal power
flow. arXiv:2212.12470, 2022.

\bibitem{rlpinn}
Optimizing energy management of smart grid using reinforcement learning aided
by surrogate models built using physics-informed neural networks.
arXiv:2510.17380, 2025.

\bibitem{saferl}
A review of safe reinforcement learning methods for modern power systems.
arXiv:2407.00304, 2024.

\bibitem{gridgeneralization}
Towards systematic generalization for power grid optimization problems.
arXiv:2605.02026, 2026.

\bibitem{pinn}
M.~Raissi, P.~Perdikaris, G.~E.~Karniadakis. Physics-informed neural
networks. \textit{Journal of Computational Physics}, 378:686--707, 2019.

\bibitem{faultsurvey}
Survey on methods for detection, classification and location of faults in
power systems using artificial intelligence. arXiv:2507.10011, 2025.

\bibitem{ntlreview}
P.~Glauner et al. The challenge of non-technical loss detection using
artificial intelligence: a survey. \textit{International Journal of
Computational Intelligence Systems}, 2017.

\bibitem{hart}
G.~W.~Hart. Nonintrusive appliance load monitoring. \textit{Proceedings of
the IEEE}, 80(12):1870--1891, 1992.

\bibitem{kelly}
J.~Kelly, W.~Knottenbelt. Neural NILM: deep neural networks applied to energy
disaggregation. \textit{ACM BuildSys}, 2015.

\bibitem{seq2point}
C.~Zhang, M.~Zhong, Z.~Wang, N.~Goddard, C.~Sutton. Sequence-to-point
learning with neural networks for non-intrusive load monitoring.
\textit{AAAI}, 2018.

\end{thebibliography}

\end{document}
